\documentclass[11pt,a4paper]{article}
\usepackage[T1]{fontenc}
\usepackage{lmodern}
\usepackage{isabelle,isabellesym}
\usepackage{amsmath,amssymb,amsthm}
\usepackage{stmaryrd}
\usepackage{wasysym}
\usepackage{textcomp}
\usepackage[margin=1in]{geometry}
\usepackage{pdfsetup}

\isabellestyle{it}
\renewcommand{\isastyletext}{\normalsize\normalfont}
\renewcommand{\isastyletxt}{\normalfont}

% bra-ket macros used in the commentary
\newcommand{\ket}[1]{|#1\rangle}
\newcommand{\bra}[1]{\langle #1|}
\newcommand{\braket}[2]{\langle #1|#2\rangle}
\newcommand{\ketbra}[2]{|#1\rangle\langle #2|}
\newcommand{\norm}[1]{\|#1\|}
\DeclareMathOperator{\Tr}{Tr}
\DeclareMathOperator{\tr}{tr}
\DeclareMathOperator{\rank}{rank}
\DeclareMathOperator{\range}{range}
\DeclareMathOperator{\spn}{span}
\DeclareMathOperator{\diag}{diag}
\DeclareMathOperator{\crk}{cr}
\newcommand{\C}{\mathbb{C}}
\newcommand{\R}{\mathbb{R}}
\newcommand{\Z}{\mathbb{Z}}
\newcommand{\Hil}{\mathcal{H}}
\newcommand{\Kil}{\mathcal{K}}

\theoremstyle{definition}
\newtheorem*{notestheorem}{Theorem (notes)}
\newtheorem*{notesdef}{Definition (notes)}
\newtheorem*{notesprop}{Proposition (notes)}
\newtheorem*{notescor}{Corollary (notes)}
\newtheorem*{noteslemma}{Lemma (notes)}
\newtheorem*{notespost}{Postulate (notes)}
\newtheorem*{notesremark}{Remark}

\begin{document}

\title{Lectures on Quantum Computing and Completely Positive Maps\\[0.5em]
  \large A formalisation in Isabelle/HOL of Vern Paulsen's Houston 2011 lecture notes}
\author{Dustin Bryant\\[0.3em]\small formal statements and proofs: the theories listed in the table of contents}
\maketitle

\begin{abstract}
This document is generated by Isabelle from the project theories.  Every
definition and theorem shown here has been checked by Isabelle/HOL (proofs are
folded).  The commentary restates Paulsen's lecture notes lecture by lecture
(Day~1 = 22~August~2011, Day~34 = 18~November~2011): definitions and results in
the notes' own language, each followed by the formal statement.  The informal
proofs of the notes are omitted; the formal proofs replace them.  Where the
formalisation deviates from the notes, the commentary says so.  The companion
file \texttt{NOTES\_AUDIT.md} lists, for every lecture, what is formalised and
what is not.
\end{abstract}

\tableofcontents
\newpage

\parindent 0pt\parskip 0.5ex

\input{session}

\part{Lectures 1--10: basic mathematics and the postulates of quantum mechanics}
\input{Basic_Math}
\input{Quantum_Computing_3_10}

\part{Lectures 11--22 and 29--32: quantum computing, density matrices, completely positive maps, error correction}
\input{Quantum_Computing11}

\part{Lecture 23: extreme points of the CP cone (Choi's Theorem 3)}
\input{Day23_ChoiExtremeUCP}

\part{Functional-analytic infrastructure (Lecture 25: Hahn--Banach)}
\input{Matrix_Typed}
\input{Essential_Functional_Analysis}
\input{Lmat_Euclidean}

\part{Lectures 26--28: separability and entanglement witnesses}
\input{Separable_Cone}
\input{Day24_Separability}

\end{document}
